\documentclass[12pt,a4paper]{article}

\usepackage[margin=1in]{geometry}
\usepackage[T1]{fontenc}
\usepackage[utf8]{inputenc}
\usepackage{lmodern}
\usepackage{enumitem}
\usepackage{hyperref}

\title{Technical Implementation (40\%)\\Jobra Hospital Management System}
\author{}
\date{\today}

\begin{document}
\maketitle

\section{Code quality and architecture}
\begin{itemize}[leftmargin=*, itemsep=0.4em]
  \item \textbf{Layered backend structure}: The backend follows a clear separation of concerns using \texttt{routes/} (HTTP endpoints), \texttt{controllers/} (request/response logic), \texttt{models/} (database operations), \texttt{middleware/} (auth and role checks), and \texttt{config/db.js} (database connection).
  \item \textbf{Role-based application design}: The system is built around three roles (admin, doctor, patient). This is reflected consistently across API route grouping (\texttt{/api/admin}, \texttt{/api/doctor}, \texttt{/api/patient}, etc.) and frontend routing (role dashboards under \texttt{/admin}, \texttt{/doctor}, \texttt{/patient}).
  \item \textbf{Frontend modularity}: The React app is organized into \texttt{pages/} (role-specific screens), \texttt{components/} (shared UI like layout/protection), \texttt{context/} (auth state), and \texttt{services/} (Axios API wrappers). This supports maintainability as features grow.
\end{itemize}

\section{Proper use of chosen technologies}
\begin{itemize}[leftmargin=*, itemsep=0.4em]
  \item \textbf{Express + REST routing}: Express is used to expose a set of JSON APIs under a consistent prefix (\texttt{/api/...}), matching the frontend base URL configuration.
  \item \textbf{MySQL integration}: The project uses \texttt{mysql2} and a schema in \texttt{database.sql}. Controllers perform CRUD operations and join queries (e.g., appointments joined with doctor/patient details).
  \item \textbf{React Router for SPA navigation}: Frontend routes are defined using \texttt{react-router-dom} with protected role-based navigation and redirects for logged-in users.
  \item \textbf{Axios service layer}: API calls are centralized via an Axios instance that injects the JWT token from \texttt{localStorage} into the \texttt{Authorization} header using an interceptor.
\end{itemize}

\section{Security implementation}
\begin{itemize}[leftmargin=*, itemsep=0.4em]
  \item \textbf{Password security}: User passwords are hashed with \texttt{bcrypt} during registration, and verified with \texttt{bcrypt.compare} during login.
  \item \textbf{Authentication}: JWT tokens are generated on login and verified by backend middleware for protected routes. Decoded tokens attach \texttt{\{id, role\}} to \texttt{req.user} for downstream authorization.
  \item \textbf{Authorization (RBAC)}: A role-check middleware enforces role permissions (e.g., admin-only pages/endpoints vs doctor/patient).
  \item \textbf{Security improvement note}: The repository includes \texttt{backend/.env} with a \texttt{JWT\_SECRET} value, but the current token generation/verification uses a hard-coded secret string (\texttt{"jobra\_secret\_key"}). For stronger security, the JWT secret should be sourced from environment variables and never hard-coded.
\end{itemize}

\section{Performance optimization}
\begin{itemize}[leftmargin=*, itemsep=0.4em]
  \item \textbf{Efficient SQL retrieval}: Some endpoints use SQL \texttt{JOIN}s to fetch relevant relational data in a single request (appointments with doctor/patient names), reducing round-trips.
  \item \textbf{Focused API design}: The backend exposes role-specific endpoints (e.g., ``my appointments'', ``doctor appointments''), which avoids sending unnecessary data to the client.
  \item \textbf{Performance improvement notes}: The current DB connector uses a single connection (\texttt{createConnection}) rather than a pool; scaling would benefit from \texttt{createPool}. Additional improvements include pagination/search for large lists (doctors, patients, appointments), and database indexing on frequently queried columns (e.g., email fields, foreign keys).
\end{itemize}

\end{document}

